\chapter{\textit{Ab initio} Quantum Many-body Methods} \label{chp:othermethods}
\epigraph{Hartree-Fock plays the same role in many-body quantum physics as a basic kitchen tool does in cooking; it can be used to prepare almost everything.}{Morten Hjorth-Jensen}
\begin{figure}[H]
	\centering
	\includegraphics[scale=0.4]{Images/example.png}
	\caption{Caption}
\end{figure}

Post Hartree-Fock methods are usually \textit{ab initio} methods where Hartree-Fock can be used as input. The term \textit{ab initio} means from first principles, implying that only physical constants are put into the methods. The Monte-Carlo methods are not considered \textit{ab initio} as nonphysical hyper parameters are required. The methods we will discuss here are the \textbf{configuration interaction} method and the \textbf{coupled cluster} method. \newpage

\section{The Hartree-Fock method}
Hartree-Fock is an algorithm for finding an approximative expression for the ground state energy of a Hamiltonian. Ever since the childhood of many-body quantum mechanics, it has been an important quantum many-body method due to its accuracy contra expensiveness. 

The method sets up a mean-field potential, often called the Hartree-Fock potential, which replaces the two-body repulsive interaction found in electronic structure calculations. Thereafter, a new basis is expanded in a given basis, and the optimal coefficients with respect to the mean-field constraint are generated. Traditionally, the method has been called the self-consistent field method, but today Hartree-Fock is mostly used. 

Since we are replacing the interaction term, we need to define a new operator to replace the Hamiltonian, known as the Hartree-Fock operator,
\begin{equation}
\hat{h}^{\text{HF}}=\hat{t} + \hat{u}_{\text{ext}} + \hat{u}^{\text{HF}}.
\end{equation}
The $\hat{u}^{\text{HF}}$ is a single-particle potential, and will later be determined by the algorithm. The operator can be used directly in the Schrödinger equation 
\begin{equation}
\hat{h}^{\text{HF}}\ket{p}=\epsilon_{\alpha}\ket{p}
\end{equation}
where the obtained energy is called the Hartree-Fock energy. The Hartree-Fock basis is found from the following expansion,
\begin{equation}
\ket{p}=\sum_{\lambda}C_{p\lambda}\ket{\lambda},
\end{equation}
where the initial basis is denoted with Greek letters and the new basis is denoted with Roman letters. $C_{p\lambda}$ are the coefficients that form an orthogonal basis. 

\subsection{Restricted Hartree-Fock}
When we calculated the reference energy above, our basis contained only one Slater determinant, more specifically the ground state. In Hartree-Fock, we still have a single Slater determinant basis, but we now construct new SPFs with the constraint of minimizing the energy. 

In general, one can change from one single-particle basis to another by a unitary transform,
\begin{equation}
\ket{p}=\sum_{\alpha}c_{p\alpha}\ket{\alpha},
\end{equation}
where we use greek letters for the old basis and roman letters for the new one. If we then insert into \eqref{eq:c_H_c}, we get a find energy formula with coefficients, $C_{p\lambda}$, that we can vary
\begin{equation}
E=\sum_{p}^N\sum_{\alpha\beta}C_{p\alpha}^*C_{p\beta}\mel{\alpha}{\hat{h}_0}{\beta}+\frac{1}{2}
\sum_{pq}^N\sum_{\alpha\beta\gamma\delta}C_{p\alpha}^*C_{q\beta}^*C_{p\gamma}C_{q\delta}\mel{\alpha\beta}{\hat{v}}{\gamma\delta}_{\text{AS}}.
\end{equation}
Further, we assume that also our new basis is orthonormal, i.e,
\begin{align}
	\braket{p}{q}&=\sum_{\alpha}c_{p\alpha}^*c_{q\alpha}\braket{\alpha}{\alpha}=\sum_{\alpha}c_{p\alpha}^*c_{q\alpha}=\delta_{pq}\\
	&\Rightarrow\sum_{\alpha}c_{p\alpha}^*c_{q\alpha}-\delta_{pq}=0\label{eq:constraint}
\end{align}

We now have a function, $E$, that we want to minimize with respect to a constraint given in equation \eqref{eq:constraint}. This is a typical situation where Lagrange Multipliers is convenient to use, which in this case can be written as
\begin{equation}
{\cal L}(\{C_{p\alpha}\})=E(\{C_{p\alpha}\})-\sum_a\varepsilon_a\Big(\sum_{\alpha}c_{p\alpha}^*c_{q\alpha}-\delta_{pq}\Big).
\end{equation}
The variation in reference energy is then find to be
\begin{equation}
\delta E=\sum_{k\alpha}\frac{\partial E}{\partial C_{k\alpha}^*}\delta C_{k\alpha}^*+\sum_{k\alpha}\frac{\partial E}{\partial C_{k\alpha}}\delta C_{k\alpha}-\sum_{k\alpha}\varepsilon_k(C_{k\alpha}\delta C_{k\alpha}^*+C_{k\alpha}^*\delta C_{k\alpha})
\label{eq:E_hf}
\end{equation}
which is zero when $E$ is minimized. Each coefficient $C_{k\alpha}$ and $C_{k\alpha}^*$ is independent, so they can be varied independently. Thus
\begin{equation}
\bigg(\frac{\partial E}{\partial C_{k\alpha}^*}-\varepsilon_kC_{k\alpha}\bigg)\delta C_{k\alpha}^*=0,
\end{equation}
which is satisfied if and only if
\begin{equation}
\frac{\partial E}{\partial C_{k\alpha}^*}-\varepsilon_kC_{k\alpha}=0\qquad\forall\,\,\, k,\alpha
\end{equation}
The first term can be derived from \eqref{eq:E_hf}, and reads
\begin{equation}
\frac{\partial E}{\partial C_{k\alpha}^*}=\sum_{\beta}C_{k\beta}\mel{\alpha}{\hat{h}_0}{\beta}+\sum_p^N\sum_{\beta\gamma\delta}C_{p\beta}^*C_{k\gamma}C_{p\delta}\mel{\alpha\beta}{\hat{v}}{\gamma\delta}_{\text{AS}}.
\end{equation}
This results in the equation
\begin{equation}
\sum_{\gamma}\hat{h}_{\alpha\gamma}^{\text{HF}}C_{k\gamma}=\varepsilon_kC_{k\gamma}
\label{eq:HF_sum}
\end{equation}
where we have defined
\begin{equation}
\hat{h}_{\alpha\gamma}^{\text{HF}}\equiv\mel{\alpha}{\hat{h}_0}{\gamma}+\sum_p^N\sum_{\beta\delta}C_{p\beta}^*C_{p\delta}\mel{\alpha\beta}{\hat{v}}{\gamma\delta}_{\text{AS}}.
\end{equation}
We recognize that \eqref{eq:HF_sum} can be written as a matrix-vector product
\begin{equation}
\hat{h}^{\text{HF}}C_k=\varepsilon_k^{\text{HF}}C_k
\end{equation}
where $C_k$ are columns in our coefficient matrix and $\varepsilon_k^{\text{HF}}$ are just the eigenvalues of $\hat{h}^{\text{HF}}$, they have no physical significance. We will use this equation to find the optimal SPFs (optimal $C_k$'s) and then find the energy from equation \eqref{eq:E_hf}.
\begin{equation}
\hat{h}^{\text{HF}}(C_k^{i+1})C_k^i=\varepsilon_k^{\text{HF}}C_k^i
\label{eq:HF_iter}
\end{equation}
Usually one initialize this with $\hat{C}=\hat{{\cal I}}$, the identity matrix.

\subsection{Unrestricted Hartree-Fock}

\section{Configuration interaction} \label{sec:ci}
The configuration interaction method is in many ways the the most intuitive method, and some will argue that it is the natural starting point. 

Often, we know the true wavefunctions $\ket{\Phi_i}$ in the external potential, but are off when interaction is added
\begin{equation}
\hat{H}_0\ket{\Phi_i}=\varepsilon_i\ket{\Phi_i},\qquad (\hat{H}_0+\hat{H}_I)\ket{\Phi_i}\neq\varepsilon_i\ket{\Phi_i}.
\end{equation}
However, the Slater determinants form a ket basis, meaning we can write out eigenstates of $\hat{H}_I$ as a linear combination of the determinants
\begin{equation}
%\begin{vmatrix}
\mqty{
	\ket{\Psi_0}&=&C_0^{(0)}\ket{\Phi_0}&+&C_1^{(0)}\ket{\Phi_1}&+&\hdots&+&C_{N-1}^{(0)}\ket{\Phi_{N-1}}\\
	\ket{\Psi_1}&=&C_0^{(1)}\ket{\Phi_0}&+&C_1^{(1)}\ket{\Phi_1}&+&\hdots&+&C_{N-1}^{(1)}\ket{\Phi_{N-1}}\\
	\ket{\Psi_2}&=&C_0^{(2)}\ket{\Phi_0}&+&C_1^{(2)}\ket{\Phi_1}&+&\hdots&+&C_{N-1}^{(2)}\ket{\Phi_{N-1}}\\
	\vdots&&\vdots&&\vdots&&&&\vdots\\
	\ket{\Psi_{N-1}}&=&C_0^{(N-1)}\ket{\Phi_0}&+&C_1^{(N-1)}\ket{\Phi_1}&+&\hdots&+&C_N^{(N-1)}\ket{\Phi_{N-1}}
}
%\end{vmatrix}
\end{equation}
such that 
\begin{equation}
\hat{H}\ket{\Psi_p}=\varepsilon_p\ket{\Psi_p}.
\end{equation}
The Hamiltonian can be rewritten as a double sum over all states using the so-called \textit{completeness relation}, 
\begin{equation}
\hat{H}=\sum_{ij}\ket{\Phi_i}\mel{\Phi_i}{\hat{H}}{\Phi_j}\bra{\Phi_j}
\end{equation}
such that the Schr\"odinger equation can be rewritten as
\begin{equation}
\begin{pmatrix}
\mel{\Phi_0}{\hat{H}}{\Phi_0} & \mel{\Phi_0}{\hat{H}}{\Phi_1} & \hdots & \mel{\Phi_0}{\hat{H}}{\Phi_{N-1}}\\
\mel{\Phi_1}{\hat{H}}{\Phi_0} & \mel{\Phi_1}{\hat{H}}{\Phi_1} & \hdots & \mel{\Phi_1}{\hat{H}}{\Phi_{N-1}}\\
\vdots & \vdots & \ddots & \vdots\\
\mel{\Phi_{N-1}}{\hat{H}}{\Phi_0} & \mel{\Phi_{N-1}}{\hat{H}}{\Phi_1} & \hdots & \mel{\Phi_{N-1}}{\hat{H}}{\Phi_{N-1}}\\
\end{pmatrix}
\begin{pmatrix}
c_0^{(p)}\\ c_1^{(p)} \\ \vdots\\ c_{N-1}^{(p)}
\end{pmatrix}
=\varepsilon_p
\begin{pmatrix}
c_0^{(p)}\\ c_1^{(p)} \\ \vdots\\ c_{N-1}^{(p)}
\end{pmatrix}
\label{eq:CIS_matrix}
\end{equation}

Until now, we have not made any assumptions, such that the equation above will give exact results when all single particle functions are included. The problem is that the matrix scales so badly, the number of Slater determinant that we need to include goes as 
\begin{equation}
N_{\text{FCI}}=\binom{N_{\text{single orbitals}}}{N_{\text{electrons}}}
\end{equation}
which is exploding. This is quite annoying since we in principle know how to solve the problems exact.

\section{Coupled Cluster} \label{subsec:cc}
The coupled cluster method is the \textit{de facto} standard wave function-based method for electronic structure calculations. \cite{paldus_beginnings_2005} The method approximates the wave function with an exponential expansion, 
\begin{equation}
\ket{\Psi_{\text{CC}}}=e^{\hat{T}}\ket{\Phi_0}
\end{equation}
where $\hatT$ is the cluster operator, entirely given by $\hatT=\hatT_1+\hatT_2 +\hatT_3+\hdots$ with
\begin{equation}
\hatT_n = \left( \frac{1}{n!}\right)^2 \sum_{abc...} \sum_{ijk...} t_{ijk...}^{abc...}a_a^\dagger a_b^\dagger a_c^\dagger \cdots a_k a_j a_i.
\end{equation}
We again want to solve the Schrödinger equation,
\begin{equation}
\hat{H}\ket{\Psi}=\hat{H}e^{\hatT}\ket{\Phi_0}=\epsilon e^{\hatT}\ket{\Phi_0},
\end{equation}
which can be simplified by multiplying with $e^{-\hatT}$ from the left. This introduces us to the \textbf{similarity transformed Hamiltonian} 
\begin{equation}
\bar{H}=e^{-\hatT}\hat{H}e^{\hatT}.
\end{equation}
If we on one hand now multiply with the reference bra on the left hand side, we easily observe that
\begin{equation}
\mel{\Phi_0}{\bar{H}}{\Phi_0}=\epsilon
\end{equation}
which is the coupled cluster energy equation. On the other hand, we can multiply with an excited bra on left hand side, and find that
\begin{equation}
\mel{\Phi_{ijk\hdots}^{abc\hdots}}{\bar{H}}{\Phi_0}=0
\end{equation}
which are the coupled cluster amplitude equations. The similarity transformed Hamiltonian can be rewritten using the Baker-Campbell-Hausdorff expansion
\begin{align}
\label{eq:BCH}
\bar{H} = \hat{H} &+ [\hat{H},\hat{T}]\notag \\
&+ \frac{1}{2}[[\hat{H},\hat{T}],\hat{T}]\notag \\
&+ \frac{1}{6}[[[\hat{H},\hat{T}],\hat{T}],\hat{T}] \\
&+ \frac{1}{24}[[[[\hat{H},\hat{T}],\hat{T}],\hat{T}],\hat{T}] \notag \\
&+ \cdots \notag
\end{align}
and we are in principle set to solve the amplitude equations with respect to the amplitudes $t_{ijk\hdots}^{abc\hdots}$ and then find the energy. The expansion is able to reproduce the true wave function exactly using a satisfying number of terms and an infinite basis. This is, of course, not possible, but even by limiting us to the first few coupled cluster operators, the results are often good compared to other methods. \cite{daniel_crawford_introduction_2007}
